\documentclass[12pt]{article}
\usepackage[a4paper]{geometry}
\usepackage{amsmath}
\usepackage{amssymb}

\title{Force Field Optimization via AWH: Gradients and Information Geometry}
\author{Alexandre Blanco González}
\date{\today}

\begin{document}

\maketitle

\section{The Accelerated Weight Histogram (AWH) Estimator}

The goal is to optimize a set of molecular mechanics force field parameters $\theta$ such that the computed free energy differences match experimental targets. The AWH method dynamically flattens the free energy landscape along a reaction coordinate $\lambda$ to ensure rigorous phase space coverage.

The AWH method updates the free energy along a $\lambda$ coordinate such that:
\begin{equation}
    F_{\lambda}^{t+1} = F_{\lambda}^t + \Delta F_{\lambda}
    \label{eq:F_update}
\end{equation}

where the update term is:
\begin{equation}
    \Delta F_{\lambda} = - \ln \frac{N \cdot \rho_{\lambda} + n \cdot w_{\lambda}\left(\theta\right)}{N \cdot \rho_{\lambda} + n \cdot \rho_{\lambda}}
    \label{eq:DeltaF}
\end{equation}

Here $N$ is the total number of samples collected during the AWH run, and $n$ is the number of samples collected since the last free energy update (effectively always 1). The term $\rho_{\lambda}$ represents the target sampling distribution. The parameters $\theta$ are the model (force field) parameters, and $w_{\lambda}\left(\theta\right)$ are the reweighting AWH weights.

These weights are defined as the normalized probability of observing a state under the bias:
\begin{equation}
    w_{\lambda}(\theta) = \frac{\rho_{\lambda} \cdot \exp\left[ F_{\lambda}^{t} - \beta_{\lambda} \cdot U_\lambda\left(\theta\right) \right]}{\sum_{k \in\Lambda} \rho_{k} \cdot \exp\left[ F_{k}^{t} - \beta_{k} \cdot U_k\left(\theta\right) \right]}
    \label{eq:weights}
\end{equation}

To map this to an optimization framework, we define a loss function based on the deviation from an experimental free energy difference $\widetilde{\Delta G}$. Let the error residual be $\mathcal{E}(\theta)$:
\begin{equation}
    \mathcal{E}(\theta) = \Delta G \left(\theta\right) - \widetilde{\Delta G}
    \label{eq:error}
\end{equation}

where $\Delta G \left(\theta\right) = F_{\lambda=1}\left(\theta\right) - F_{\lambda = 0} \left(\theta\right)$.

\section{First-Order Gradients of the Estimator}

To minimize the loss $\mathcal{L}(\mathcal{E})$, we compute the gradient with respect to the parameters $\theta$. By the chain rule:
\begin{equation}
    \nabla_{\theta}\mathcal{L} = \frac{d\mathcal{L}}{d\mathcal{E}} \cdot \nabla_{\theta} \mathcal{E}(\theta) = \frac{d\mathcal{L}}{d\mathcal{E}} \cdot \left[ \nabla_{\theta} F_{\lambda=1}\left(\theta\right) - \nabla_{\theta} F_{\lambda = 0} \left(\theta\right) \right]
    \label{eq:chain_rule}
\end{equation}

We focus on deriving the gradient of the free energy estimator $F_{\lambda}(\theta)$. Assuming the history $F^t$ is fixed (detached from the current graph):
\begin{equation}
    \nabla_{\theta} F_{\lambda}\left(\theta\right) = \nabla_{\theta} \left[ F_{\lambda}^{t} + \Delta F_{\lambda}\left( \theta \right) \right] = \nabla_{\theta} \Delta F_{\lambda}\left(\theta\right)
    \label{eq:grad_F_step1}
\end{equation}

Substituting $\Delta F_{\lambda}$ (Equation \ref{eq:DeltaF}):
\begin{equation}
    \nabla_{\theta} \Delta F_{\lambda}\left(\theta\right) = \nabla_{\theta} \left[ - \ln \left( \frac{N \cdot \rho_{\lambda} + n \cdot w_{\lambda}\left(\theta\right)}{N \cdot \rho_{\lambda} + n \cdot \rho_{\lambda}} \right) \right]
    \label{eq:grad_F_step2}
\end{equation}

The denominator inside the logarithm is constant w.r.t $\theta$. Therefore, we only differentiate the numerator term, let us call it $u(\theta) = N \rho_{\lambda} + n w_{\lambda}(\theta)$:
\begin{equation}
    \nabla_{\theta} \left[ - \ln \left( u(\theta) \right) \right] = \frac{-1}{u(\theta)} \cdot \nabla_{\theta} u(\theta)
    \label{eq:grad_log}
\end{equation}

The gradient of the argument $u(\theta)$ is:
\begin{equation}
    \nabla_{\theta} u\left( \theta \right) = n \cdot \nabla_{\theta} w_{\lambda}\left( \theta \right)
    \label{eq:grad_u}
\end{equation}

Combining Equations \ref{eq:grad_F_step2} through \ref{eq:grad_u} yields:
\begin{equation}
    \nabla_{\theta} F_{\lambda} \left( \theta \right) = \frac{-n \cdot \nabla_{\theta} w_{\lambda}\left( \theta \right)}{N \cdot \rho_{\lambda} + n \cdot w_{\lambda}(\theta)}
    \label{eq:grad_F_u}
\end{equation}

To find $\nabla_{\theta} w_{\lambda}$, we apply the quotient rule to the weight definition (Equation \ref{eq:weights}). While direct expansion yields a complex sum of exponentials, the expression simplifies significantly by identifying the expectation values. Let $Z = \sum_{k \in \Lambda} \rho_k \exp(F_k^t - \beta_k U_k)$ be the partition function denominator. We can express the gradient of $w_{\lambda}$ as:
\begin{equation}
    \nabla_{\theta} w_{\lambda} = \frac{(\nabla_{\theta}\mathcal{N}) Z - \mathcal{N} (\nabla_{\theta} Z)}{Z^{2}} = \frac{\nabla_{\theta}\mathcal{N}}{Z} - w_{\lambda} \frac{\nabla_{\theta} Z}{Z}
    \label{eq:grad_w_quotient}
\end{equation}

Substituting the derivatives of the exponential terms:
\begin{equation}
    \frac{\nabla_{\theta}\mathcal{N}}{Z} = - \beta_{\lambda} \cdot \nabla_{\theta} U_{\lambda} \cdot w_{\lambda}
    \label{eq:grad_N}
\end{equation}

\begin{align}
    \begin{split}
        \frac{\nabla_{\theta} Z}{Z} &= \sum_{k \in \Lambda} - \beta_k \cdot \nabla_{\theta} U_k \frac{\rho_k \cdot \exp\left(F_{k}^t - \beta_{k} \cdot U_{k}\left( \theta \right) \right)}{Z}\\ 
        &= - \sum_{k \in \Lambda} w_k \cdot \beta_k \cdot \nabla_{\theta} U_k\\
        &= - \langle \beta \nabla_{\theta} U \rangle_w
    \end{split}
    \label{eq:grad_Z}
\end{align}

Combining Equations \ref{eq:grad_N} and \ref{eq:grad_Z} into Equation \ref{eq:grad_w_quotient} gives the compact form:
\begin{equation}
    \nabla_{\theta} w_{\lambda} = w_{\lambda} \left[ \langle \beta \nabla_{\theta} U \rangle_w - \beta_{\lambda} \nabla_{\theta} U_{\lambda} \right]
    \label{eq:grad_w_compact}
\end{equation}

where $\langle \beta \nabla_{\theta} U \rangle_w$ is the ensemble average of the energy gradients weighted by the current AWH weights. We substitute Equation \ref{eq:grad_w_compact} into the gradient of the free energy estimator (Equation \ref{eq:grad_F_u}) to obtain:
\begin{equation}
    \nabla_{\theta} F_{\lambda} \left( \theta \right) = \frac{n \cdot w_{\lambda}(\theta)}{N \cdot \rho_{\lambda} + n \cdot w_{\lambda}(\theta)} \left[ \beta_{\lambda} \nabla_{\theta} U_{\lambda} - \langle \beta \nabla_{\theta} U \rangle_w \right]
    \label{eq:grad_F_final}
\end{equation}

We define the prefactor $A(\lambda)$ as the effective update rate for a given window:
\begin{equation}
    A(\lambda) = \frac{n \cdot w_{\lambda}(\theta)}{N \cdot \rho_{\lambda} + n \cdot w_{\lambda}(\theta)}
    \label{eq:A_lambda}
\end{equation}

Substituting Equation \ref{eq:grad_F_final} into the chain rule (Equation \ref{eq:chain_rule}) and grouping the ensemble average terms, we arrive at the simplified gradient of the loss:
\begin{equation}
    \nabla_{\theta}\mathcal{L} = \frac{\partial \mathcal{L}}{\partial \mathcal{E}} \left[ A(1) \beta_1 \nabla_{\theta} U_1 - A(0) \beta_0 \nabla_{\theta} U_0 - \left( A(1) - A(0) \right) \langle \beta \nabla_{\theta} U \rangle_w \right]
    \label{eq:grad_loss_final}
\end{equation}

This form highlights that the global ensemble average $\langle \beta \nabla_{\theta} U \rangle_w$ in Equation \ref{eq:grad_loss_final} only contributes to the gradient if there is an imbalance in the effective sampling rate between the end states $\lambda=1$ and $\lambda=0$.

\section{Loss Function Geometries}

The choice of loss function $\mathcal{L}(\mathcal{E})$ determines the pre-factor governing how strongly the optimizer reacts to deviations from experimental data.

The L1 loss is robust to outliers, which is useful when initial free energy estimates are very poor.
\begin{equation}
    \mathcal{L}_{L1} = | \mathcal{E} |, \quad \frac{d\mathcal{L}_{L1}}{d\mathcal{E}} = \mathrm{sign}(\mathcal{E})
    \label{eq:loss_L1}
\end{equation}

This provides a constant gradient magnitude regardless of how far the prediction is from the target, but is non-differentiable at $\mathcal{E}=0$.

The L2 loss penalizes large errors more heavily and provides a smooth gradient that approaches zero as the prediction converges.
\begin{equation}
    \mathcal{L}_{L2} = \mathcal{E}^2, \quad \frac{d\mathcal{L}_{L2}}{d\mathcal{E}} = 2 \mathcal{E}
    \label{eq:loss_L2}
\end{equation}

Note that if defined as $\frac{1}{2}\mathcal{E}^2$, the derivative simplifies to just $\mathcal{E}$. The L2 loss is generally preferred for fine-tuning but can be unstable if gradients are extremely large initially.

The Huber loss combines the best properties of L1 and L2. It is quadratic for small errors (stable convergence) and linear for large errors (robust to outliers). Defined with a threshold $\delta$:
\begin{equation}
    \mathcal{L}_{\delta} = \begin{cases}
    \frac{1}{2}\mathcal{E}^2 & \text{for } |\mathcal{E}| \le \delta, \\
    \delta (|\mathcal{E}| - \frac{1}{2}\delta), & \text{otherwise.}
    \end{cases}
    \label{eq:loss_huber}
\end{equation}

The derivative is:
\begin{equation}
    \frac{d\mathcal{L}_{\delta}}{d\mathcal{E}} = \begin{cases}
    \mathcal{E} & \text{for } |\mathcal{E}| \le \delta, \\
    \delta \cdot \mathrm{sign}(\mathcal{E}), & \text{otherwise.}
    \end{cases}
    \label{eq:deriv_huber}
\end{equation}

This ensures that the gradients used to update the force field parameters never exceed magnitude $\delta$, preventing numerical instability during the early stages of optimization.

\section{Information Theory and the Fisher Information Matrix}

While the Huber loss clips the magnitude of the gradient, it does not correct its direction in the underlying parameter space. Standard gradient descent assumes the parameter space $\theta$ is Euclidean, meaning a step of magnitude $\Delta$ in one parameter is geometrically equivalent to a step of magnitude $\Delta$ in another. In molecular mechanics, this assumption fails severely. A perturbation of $0.05$ units in a Lennard-Jones collision radius ($\sigma$) induces massive steric clashes and fundamentally alters the sampled ensemble. Conversely, the same $0.05$ perturbation in a soft dihedral force constant barely affects the total energy, leaving the sampled ensemble largely unchanged. 

Instead of an Euclidean space, the parameters $\theta$ define coordinates on a \textit{statistical manifold}. To guarantee valid ensemble reweighting, the optimization step must be constrained by the geometry of this manifold. 

In the proposed algorithmic workflow, AWH is run until the free energy estimates $F_\lambda$ converge. The bias is then fixed, and the system explores an extended ensemble defined jointly over the atomic coordinates $\mathbf{x}$ and the alchemical states $\lambda$. Each point on our statistical manifold is therefore not a simple canonical ensemble, but the joint probability distribution of this extended AWH ensemble:
\begin{equation}
    P_{\theta}(\mathbf{x}, \lambda) = \frac{\rho_\lambda \exp\left(F_\lambda - \beta_\lambda U_\lambda(\theta, \mathbf{x})\right)}{\mathcal{Z}(\theta)}
    \label{eq:P_theta}
\end{equation}

where $\mathcal{Z}(\theta)$ is the global partition function of the extended ensemble:
\begin{equation}
    \mathcal{Z}(\theta) = \sum_{k \in \Lambda} \int \rho_k \exp\left(F_k - \beta_k U_k(\theta, \mathbf{x})\right) d\mathbf{x}
    \label{eq:Z_theta_global}
\end{equation}

Note that averaging an observable over this joint distribution $P_{\theta}(\mathbf{x}, \lambda)$ corresponds exactly to the AWH-weighted ensemble average $\langle \dots \rangle_w$ utilized in our gradient estimators.

To optimize safely, we must measure the statistical distance between the extended ensemble generated at the current epoch $P_{\theta}$ and the proposed distribution $P_{\theta+d\theta}$, ensuring that both have sufficient phase-space overlap. In information theory, this distance is quantified by the Kullback-Leibler (KL) divergence:
\begin{align}
    \begin{split}
        D_{KL}(P_{\theta} || P_{\theta+d\theta}) &= \sum_{\lambda \in \Lambda} \int P_{\theta}(\mathbf{x}, \lambda) \ln \frac{P_{\theta}(\mathbf{x}, \lambda)}{P_{\theta+d\theta}(\mathbf{x}, \lambda)} d\mathbf{x} \\
        &= \left \langle \ln{\frac{P_{\theta}(\mathbf{x}, \lambda)}{P_{\theta+d\theta}(\mathbf{x}, \lambda)}} \right\rangle_{w}
    \end{split}
    \label{eq:KL_div}
\end{align}

To relate small parameter updates $d\theta$ to this statistical distance, we perform a second-order Taylor expansion of $\ln P_{\theta+d\theta}$ around $\theta$:
\begin{equation}
    \ln P_{\theta+d\theta} \approx \ln P_{\theta} + d\theta^T \nabla_{\theta} \ln P_{\theta} + \frac{1}{2} d\theta^T \nabla^2_{\theta} \ln P_{\theta} d\theta
    \label{eq:taylor_expansion}
\end{equation}

Substituting this expansion back into the KL divergence (Equation \ref{eq:KL_div}) yields:
\begin{equation}
    D_{KL}(P_{\theta} || P_{\theta+d\theta}) \approx -\left\langle d\theta^T \nabla_{\theta} \ln P_{\theta} + \frac{1}{2} d\theta^T \nabla^2_{\theta} \ln P_{\theta} d\theta \right\rangle_w
    \label{eq:KL_taylor}
\end{equation}

The expectation of the first-order score function $\nabla_{\theta} \ln P_{\theta}$ strictly vanishes:
\begin{equation}
    \left\langle \nabla_{\theta} \ln P_{\theta} \right\rangle_w = \sum_{\lambda} \int P_{\theta} \frac{\nabla_{\theta} P_{\theta}}{P_{\theta}} d\mathbf{x} = \nabla_{\theta} \left( \sum_{\lambda} \int P_{\theta} d\mathbf{x} \right) = \nabla_{\theta} (1) = 0
    \label{eq:score_vanishes}
\end{equation}

This leaves the second-order term:
\begin{equation}
    D_{KL}(P_{\theta} || P_{\theta+d\theta}) \approx - \frac{1}{2} d\theta^T \left\langle \nabla^2_{\theta} \ln P_{\theta} \right\rangle_w d\theta
    \label{eq:KL_second_order}
\end{equation}

The matrix $-\langle \nabla^2_{\theta} \ln P_{\theta} \rangle_w$ is the Fisher Information Matrix (FIM), $\mathcal{I}(\theta)$. Under standard regularity conditions, the negative expectation of the Hessian of the log-likelihood is mathematically equivalent to the covariance of the score functions:
\begin{equation}
    \mathcal{I}_{ij}(\theta) = \left\langle \left( \nabla_{\theta_i} \ln P_{\theta} \right) \left( \nabla_{\theta_j} \ln P_{\theta} \right) \right\rangle_w
    \label{eq:FIM_def}
\end{equation}

Thus, the FIM acts as the metric tensor for the statistical manifold, explicitly translating parameter updates into statistical distance:
\begin{equation}
    D_{KL}(P_{\theta} || P_{\theta+d\theta}) \approx \frac{1}{2} d\theta^T \mathcal{I}(\theta) d\theta
    \label{eq:KL_FIM}
\end{equation}

We can now analytically derive the FIM for the AWH extended ensemble. First, we compute the score function $\nabla_{\theta} \ln P_{\theta}(\mathbf{x}, \lambda)$:
\begin{align}
    \nabla_{\theta} \ln P_{\theta}(\mathbf{x}, \lambda) &= \nabla_{\theta} \left[ \ln \rho_\lambda + F_\lambda - \beta_\lambda U_\lambda(\theta, \mathbf{x}) - \ln \mathcal{Z}(\theta) \right] \nonumber \\
    &= -\beta_\lambda \nabla_{\theta} U_\lambda(\theta, \mathbf{x}) - \frac{\nabla_{\theta} \mathcal{Z}(\theta)}{\mathcal{Z}(\theta)}
    \label{eq:score_deriv_1}
\end{align}

Because $\mathcal{Z}(\theta)$ is exactly proportional to the partition function denominator $Z$ utilized in Equation \ref{eq:grad_w_quotient}, we have already established its derivative in Equation \ref{eq:grad_Z}:
\begin{equation}
    \frac{\nabla_{\theta} \mathcal{Z}(\theta)}{\mathcal{Z}(\theta)} = -\langle \beta \nabla_{\theta} U \rangle_{w}
    \label{eq:grad_Z_global}
\end{equation}

Therefore, the score function of the extended ensemble is precisely the deviation of the local energy gradient from its global AWH-weighted ensemble average:
\begin{equation}
    \nabla_{\theta} \ln P_{\theta}(\mathbf{x}, \lambda) = -\beta_\lambda \nabla_{\theta} U_\lambda(\theta, \mathbf{x}) + \langle \beta \nabla_{\theta} U \rangle_{w}
    \label{eq:score_final}
\end{equation}

Substituting this score function into the FIM definition (Equation \ref{eq:FIM_def}) yields exactly the covariance of the energy gradients evaluated over the extended AWH ensemble:
\begin{align}
    \begin{split}
        \mathcal{I}_{ij}(\theta) &= \left\langle \left[ -\beta_\lambda \nabla_{\theta_i} U_\lambda + \langle \beta \nabla_{\theta_i} U \rangle_w \right] \left[ -\beta_\lambda \nabla_{\theta_j} U_\lambda + \langle \beta \nabla_{\theta_j} U \rangle_w \right] \right\rangle_w \\
        &= \langle \beta^2 \nabla_{\theta_i} U \nabla_{\theta_j} U \rangle_w - \langle \beta \nabla_{\theta_i} U \rangle_w \langle \beta \nabla_{\theta_j} U \rangle_w \\
        &= \text{Cov}_w(\beta \nabla_{\theta_i} U, \beta \nabla_{\theta_j} U)
    \end{split}
    \label{eq:FIM_cov}
\end{align}

This derivation reveals why optimizing the loss via standard gradient descent is problematic for reweighting schemes. If the optimizer alters a sensitive parameter based solely on the loss gradient $\nabla_{\theta} \mathcal{L}$, the resulting distribution $P_{\theta+d\theta}$ may have zero overlap with the proposal distribution $P_{\theta}$. If the distributions do not overlap, the Effective Sample Size (ESS) collapses, weight degeneracy occurs, and the ensemble average estimators $\langle \beta \nabla_{\theta} U \rangle_w$ become pure noise.

To guarantee valid ensemble reweighting across optimization epochs, we must scale the gradient step by the statistical curvature. Updating parameters using the natural gradient $\Delta \theta = -\eta \mathcal{I}^{-1}(\theta) \nabla_{\theta} \mathcal{L}$ directly bounds the KL divergence of the step. By penalizing updates to sensitive parameters (large variance in $\nabla_{\theta} U$) and accelerating updates to stiff parameters (small variance in $\nabla_{\theta} U$), the natural gradient prevents parameter updates from stepping outside the phase-space overlap of the generated AWH trajectory.

\section{Algorithmic Implementation}

To integrate the theoretical derivations into a functional machine learning optimization pipeline, we partition the algorithm into three distinct phases: reference ensemble generation, natural gradient optimization, and resimulation. The approach dynamically reweights parameter trajectories until the statistical overlap deteriorates.

\subsection{Phase 1: Reference Ensemble Generation}

This phase constructs the empirical dataset that acts as the support for subsequent optimization steps. 
\begin{enumerate}
    \item \textbf{AWH Convergence:} Initialize the molecular dynamics simulation with the baseline force field parameters $\theta_0$. Execute the AWH enhanced sampling protocol until the free energy estimates $F_\lambda$ reach convergence criteria across all alchemical states $\lambda \in \Lambda$.
    \item \textbf{Extended Ensemble Production:} Fix the bias $F_\lambda$ to halt its adaptation. Perform a production run exploring the extended ensemble $(\mathbf{x}, \lambda)$, utilizing a Gibbs sampler to transition between $\lambda$ states.
    \item \textbf{Trajectory Logging:} Record $M$ frames from the production run. For each frame index $k \in \{1, \dots, M\}$, store the atomic coordinates $\mathbf{x}_k$ and the instantaneously active alchemical state $\lambda_k$. 
    \item \textbf{Baseline Energy Evaluation:} Compute the dimensionless potential energy for all $M$ frames using the baseline parameter set $\theta_0$:
    \begin{equation}
        u_k^{(0)} = \beta_{\lambda_k} U_{\lambda_k}(\theta_0, \mathbf{x}_k)
        \label{eq:baseline_energy}
    \end{equation}
\end{enumerate}

\subsection{Phase 2: Natural Gradient Optimization Loop}

This loop exploits the generated $M$ frames to take multiple optimization steps ($\theta_0 \to \theta_1 \dots \to \theta_n$). During this phase, any extended ensemble expectation value $\langle \mathcal{O} \rangle_w$ is evaluated as a discretely weighted empirical mean. For the current epoch $n$:

\begin{enumerate}
    \item \textbf{Forward Pass and Gradients:} For each frame $k$, evaluate the updated dimensionless energy $u_k^{(n)}$ and compute the gradient vector of the energy with respect to the parameters:
    \begin{align}
        u_k^{(n)} &= \beta_{\lambda_k} U_{\lambda_k}(\theta_n, \mathbf{x}_k) \label{eq:epoch_energy} \\
        \mathbf{s}_k^{(n)} &= \nabla_\theta \left[ \beta_{\lambda_k} U_{\lambda_k}(\theta_n, \mathbf{x}_k) \right] \label{eq:epoch_gradient}
    \end{align}
    \item \textbf{Parameter Reweighting:} Calculate the unnormalized importance sampling weight for each frame relative to the reference ensemble sampled in Phase 1:
    \begin{equation}
        W_k^{(n)} = \exp\left( - \left[ u_k^{(n)} - u_k^{(0)} \right] \right)
        \label{eq:unnorm_weight}
    \end{equation}
    Normalize these weights to sum to unity:
    \begin{equation}
        w_k^{(n)} = \frac{W_k^{(n)}}{\sum_{j=1}^M W_j^{(n)}}
        \label{eq:norm_weight}
    \end{equation}
    \item \textbf{Phase Space Overlap Verification:} Calculate the Effective Sample Size (ESS) to quantify the statistical overlap between the target and reference distributions:
    \begin{equation}
        \text{ESS}^{(n)} = \frac{1}{\sum_{k=1}^M \left(w_k^{(n)}\right)^2}
        \label{eq:ess}
    \end{equation}
    If $\text{ESS}^{(n)} < \text{ESS}_{\text{threshold}}$ (where the threshold is typically a predefined fraction of $M$, e.g., $0.1M$), the reference frames no longer possess the statistical support necessary to optimize $\theta_n$. The optimization loop is terminated, and the algorithm proceeds to Phase 3.
    \item \textbf{Constructing the Preconditioner:} Compute the weighted mean of the energy gradient vectors:
    \begin{equation}
        \overline{\mathbf{s}}^{(n)} = \sum_{k=1}^M w_k^{(n)} \mathbf{s}_k^{(n)}
        \label{eq:mean_gradient}
    \end{equation}
    Note that $\overline{\mathbf{s}}^{(n)}$ is exactly the empirical estimator for the term $\langle \beta \nabla_{\theta} U \rangle_w$ defined in Equation \ref{eq:grad_loss_final}. Evaluate the Fisher Information Matrix (FIM) as the empirical covariance matrix over the reweighted ensemble, substituting the terms derived in Equation \ref{eq:FIM_cov}:
    \begin{equation}
        \mathcal{I}(\theta_n) = \sum_{k=1}^M w_k^{(n)} \left( \mathbf{s}_k^{(n)} {\mathbf{s}_k^{(n)}}^T \right) - \overline{\mathbf{s}}^{(n)} {\overline{\mathbf{s}}^{(n)}}^T
        \label{eq:empirical_fim}
    \end{equation}
    \item \textbf{Loss Gradient Evaluation:} Evaluate the target loss gradient $\mathbf{g}_n = \nabla_\theta \mathcal{L}(\theta_n)$ using the analytical expression defined in Equation \ref{eq:grad_loss_final}. The endpoint potentials $U_1$ and $U_0$ and their gradients are evaluated exclusively over the frames localized at $\lambda=1$ and $\lambda=0$.
    \item \textbf{Parameter Update:} Solve the linear system dictated by the natural gradient descent rule to find the update direction $\mathbf{d}_n$:
    \begin{equation}
        \mathcal{I}(\theta_n) \mathbf{d}_n = \mathbf{g}_n
        \label{eq:linear_system}
    \end{equation}
    Update the parameter vector using a learning rate $\eta$:
    \begin{equation}
        \theta_{n+1} = \theta_n - \eta \mathbf{d}_n
        \label{eq:natural_grad_update}
    \end{equation}
    Following the update, increment $n$ and return to Step 1 of Phase 2.
\end{enumerate}

\subsection{Phase 3: Resimulation}

When the ESS collapse triggers an exit from Phase 2, the current force field $\theta_n$ has diverged outside the trust region of the original simulation.
\begin{enumerate}
    \item \textbf{System Reinitialization:} Instantiate a new molecular dynamics engine using the recently optimized parameter set $\theta_n$.
    \item \textbf{Bias Injection:} To circumvent the computationally expensive initial flattening phase, extract the converged free energy bias vector $F_\lambda$ and the target distributions $\rho_\lambda$ from the Phase 1 simulation and explicitly inject them into the new simulation state.
    \item \textbf{Restart:} Return to Step 2 of Phase 1 to execute a fresh production simulation, generating a new trajectory of $M$ frames to act as the reference ensemble for the next sequence of optimization epochs.
\end{enumerate}

\end{document}